% 网络单纯形法的一轮：可行树基 -> 求树流 -> 求势 -> 定价 -> (无违反: 停止) / 入弧与基本环 -> 推送到碰界 -> 出弧离树 -> 新树基。
\begin{tikzpicture}[>=Stealth,font=\fontsize{8}{10}\selectfont,
  box/.style={draw=FigureStroke,fill=FigureBlueFill,rounded corners=2pt,line width=.7pt,align=center,inner sep=3pt,minimum height=9mm,text width=21mm},
  stop/.style={box,fill=FigureGreenFill},
  ar/.style={->,draw=FigureStroke,line width=.8pt},
  lab/.style={text=FigureMuted,font=\FigureNoteFont\fontsize{7.5}{9}\selectfont}]
\node[box] (b1) at (0,0) {可行树基\\$(T,L,U)$};
\node[box] (b2) at (2.95,0) {从叶子起\\用守恒式求树弧流量};
\node[box] (b3) at (5.9,0) {令树弧$\bar c=0$\\沿树求节点势$\pi$};
\node[box] (b4) at (8.85,0) {定价：找符号\\不对的非树弧};
\node[stop] (b5) at (11.8,0) {停止：流与势\\构成最优性证书};
\node[box] (b6) at (8.85,-1.8) {入弧＋树上路径\\＝基本环};
\node[box] (b7) at (5.9,-1.8) {沿环推送，\\直到某条弧碰界};
\node[box] (b8) at (2.95,-1.8) {碰界的树弧出树、\\入弧进树（或\\入弧翻转界状态）};
\draw[ar] (b1)--(b2); \draw[ar] (b2)--(b3); \draw[ar] (b3)--(b4);
\draw[ar] (b4)--node[lab,above]{没有}(b5);
\draw[ar] (b4)--node[lab,right]{有}(b6);
\draw[ar] (b6)--(b7); \draw[ar] (b7)--(b8);
\draw[ar] (b8)-|(b1);
\end{tikzpicture}
